%% ============================================================
%%  AlzDetect — Journal manuscript (IEEEtran)
%%  Target: IEEE Access / SN Computer Science (adjust template per venue)
%%
%%  HOW TO USE
%%   1. Upload this folder to Overleaf (New Project -> Upload).
%%   2. Set compiler to pdfLaTeX, main document = main.tex.
%%   3. Search for "TODO" — every one marks a number to verify or a
%%      figure/experiment to insert before submission.
%%
%%  INTEGRITY NOTE: the results below use the most detailed table from the
%%  capstone (63% -> 86%). Your abstract (56% -> 92%) and GitHub README
%%  (94.2% -> 96.1%) report DIFFERENT numbers. Pick ONE clean reproducible
%%  run and make every number in this paper match it.
%% ============================================================
\documentclass[journal]{IEEEtran}

\usepackage{graphicx}
\usepackage{amsmath}
\usepackage{booktabs}
\usepackage{array}
\usepackage{url}
\usepackage{hyperref}

\begin{document}

\title{Fuzzy Logic--Based Adaptive Resampling for Class-Imbalanced
Alzheimer's Disease Staging from MRI using MobileNetV2}

\author{Aditya Jaiswal, Nilabh Shivam, Parth Vaishnav, Prabhat Tiwari,
and Stiena Wilson Anthony% TODO: confirm author order; add supervisor as corresponding author
\thanks{The authors are with the School of Computing Science and Engineering,
VIT Bhopal University, Madhya Pradesh, India (e-mail: TODO).}}

\markboth{Preprint --- Draft for Submission}{}
\maketitle

\begin{abstract}
Early staging of Alzheimer's disease (AD) from magnetic resonance imaging (MRI)
is hampered by severe class imbalance: non-demented scans vastly outnumber
moderate-dementia cases, biasing convolutional neural networks (CNNs) toward the
majority class and collapsing recall on the rarest, most clinically important
stages. We present a fuzzy logic--based adaptive resampling scheme that sets a
per-class oversampling factor as a function of each class's imbalance ratio
using a triangular-membership Mamdani controller, and couples it with a
transfer-learned MobileNetV2 classifier for four-way AD staging (Non-Demented,
Very Mild, Mild, Moderate). On the Kaggle Alzheimer MRI dataset, the adaptive
scheme raised overall accuracy from \textbf{63\% to 86\%}% TODO: verify from single clean run
and macro-averaged F1 from \textbf{0.49 to 0.85}, with moderate-dementia recall
improving from \textbf{0\% to 90\%}. We compare the fuzzy controller against
standard imbalance baselines (class weighting, random oversampling, SMOTE, and
focal loss)% TODO: run these baselines — required for a journal
and report 5-fold cross-validated results. We discuss the method's limitations,
including dataset-level leakage risks and the modest novelty of the controller
relative to proportional oversampling.
\end{abstract}

\begin{IEEEkeywords}
Alzheimer's disease, fuzzy logic, class imbalance, MobileNetV2, transfer
learning, medical image classification, MRI.
\end{IEEEkeywords}

\section{Introduction}
\IEEEPARstart{A}{lzheimer's} disease is the leading cause of dementia and a
progressive neurodegenerative disorder for which early, accurate staging is
critical to timely intervention. Deep CNNs applied to structural MRI have shown
strong performance on AD detection, but their accuracy degrades sharply when the
training data are class-imbalanced---a near-universal property of clinical
neuroimaging cohorts, where healthy and very-mild cases dominate. A model trained
naively on such data tends to maximise overall accuracy by predicting the
majority class, yielding near-zero recall on moderate dementia despite this being
the most consequential error to avoid.

In our baseline experiments, an unmodified MobileNetV2 trained on the raw
imbalanced data achieved only \textbf{63\%}% TODO: verify
overall accuracy and \textbf{0\%} recall on moderate dementia. This paper asks
whether a principled, \emph{adaptive} resampling policy---one that allocates more
synthetic samples to classes in proportion to how imbalanced they are---can
recover performance on the rare stages without the manual tuning of fixed-ratio
oversampling. We formulate this policy as a small fuzzy inference system.

\textbf{Contributions.}
(1) A fuzzy Mamdani controller that maps a class's imbalance ratio to an
oversampling factor via triangular membership functions and three IF--THEN rules.
(2) Integration of this controller as a preprocessing stage for a transfer-learned
MobileNetV2 four-way AD stager.
(3) An empirical comparison against standard imbalance-handling baselines, with
per-class metrics and cross-validation.% TODO: complete baselines + CV

\section{Related Work}
Hybrid fuzzy--deep-learning systems have been explored across medical and
non-medical imaging. Fuzzy pooling within CNNs has improved lung-ultrasound
classification with explainable outputs~\cite{ref2}. Surveys of fuzzy deep
learning for AD highlight fuzzy logic's role in managing annotation uncertainty
and improving interpretability~\cite{ref3,ref8}. Fuzzy--CNN hybrids have been
applied to edge-detected AD MRI~\cite{ref1}, melanoma diagnosis via fuzzy
correlation maps~\cite{ref10}, three-way opinion classification~\cite{ref4},
handwritten-digit recognition~\cite{ref6}, and general image
classification~\cite{ref7}. Transfer learning with CNN backbones is a standard
route to high accuracy on limited data~\cite{ref9}, and broad reviews catalogue
machine-learning approaches to AD detection~\cite{ref5}. In contrast to prior
work that embeds fuzzy operations \emph{inside} the network, we keep the CNN
unchanged and apply fuzzy reasoning at the \emph{data-balancing} stage, which
isolates the contribution of adaptive resampling and keeps the pipeline simple
and reproducible.

\section{Methodology}
\subsection{Overview}
The pipeline has three stages: (i) standardised MRI preprocessing, (ii) a fuzzy
adaptive-resampling controller that balances the training set, and (iii) a
transfer-learned MobileNetV2 classifier. Fig.~\ref{fig:arch} shows the
architecture.% TODO: insert architecture figure as arch.png
\begin{figure}[t]
  \centering
  % \includegraphics[width=\columnwidth]{arch.png}  % TODO: add figure file
  \fbox{\parbox[c][3.2cm][c]{0.9\columnwidth}{\centering TODO: architecture figure
  (Input MRI $\rightarrow$ Preprocess $\rightarrow$ Fuzzy resampling $\rightarrow$
  MobileNetV2 $\rightarrow$ 4-class softmax)}}
  \caption{Proposed fuzzy-resampling + MobileNetV2 pipeline.}
  \label{fig:arch}
\end{figure}

\subsection{Preprocessing}
MRI slices are resized to $128\times128$, converted to three channels, and
intensity-normalised to $[0,1]$. Labels are one-hot encoded over the four stages.

\subsection{Fuzzy Adaptive-Resampling Controller}
For each class $c$ with sample count $n_c$, we define an imbalance score
\begin{equation}
  I_c = 1 - \frac{n_c}{\max_k n_k},
\end{equation}
so the largest class has $I_c = 0$. A single-input, single-output Mamdani fuzzy
system maps $I_c$ to a resampling factor $r_c \in [0,1]$. The antecedent
\textit{imbalance} and consequent \textit{resample\_factor} each use three
\emph{static} triangular membership functions (low, medium, high), governed by
three rules:
\begin{itemize}
  \item R1: \textbf{IF} imbalance is low \textbf{THEN} resample is low;
  \item R2: \textbf{IF} imbalance is medium \textbf{THEN} resample is medium;
  \item R3: \textbf{IF} imbalance is high \textbf{THEN} resample is high.
\end{itemize}
Centroid defuzzification yields a crisp $r_c$, and the target size for class $c$
is $n_c^{\ast} = n_c + r_c \cdot \max_k n_k$. Minority classes are augmented up to
$n_c^{\ast}$ using label-preserving transformations (rotation $\pm15^\circ$,
width/height shift $\pm10\%$, horizontal flip, brightness $\pm20\%$).

\noindent\textit{Remark.} The membership functions and rules are fixed, not
learned. We deliberately do not claim attention gating or training-time
adaptation; the contribution is an interpretable, parameter-light alternative to
fixed-ratio oversampling.% Honesty guard — keeps claims aligned with the code.

\subsection{Classifier}
We use MobileNetV2 (ImageNet-pretrained) with all but the last 20 layers frozen,
followed by global average pooling, a 128-unit ReLU dense layer, dropout (0.3),
and a 4-way softmax. Training uses Adam (lr $=10^{-4}$), categorical
cross-entropy, batch size 32, with early stopping on validation loss.

\section{Experimental Setup}
\subsection{Dataset}
We use the Kaggle Alzheimer MRI dataset (four classes: Non-Demented, Very Mild,
Mild, Moderate Demented).% TODO: cite exact dataset + version + class counts
\textbf{TODO:} report exact per-class counts and, critically, whether the split
is \emph{patient-level} or \emph{image-level}. Image-level splitting on this
dataset risks leakage (augmented slices of one subject across train/test) and
must be addressed or explicitly disclosed.

\subsection{Protocol}
\textbf{TODO: reconcile.} The capstone reports both an 80/20 and a 50/20/30
split. Fix one. For a journal we recommend stratified \textbf{5-fold
cross-validation}, reporting mean~$\pm$~std. Metrics: overall accuracy, and
per-class precision, recall, and F1, plus macro-averaged F1 and confusion
matrices.

\section{Results}
\subsection{Effect of Fuzzy Adaptive Resampling}
Table~\ref{tab:main} summarises the improvement over the imbalanced baseline.
% TODO: all numbers below must come from ONE clean reproducible run.
\begin{table}[t]
\caption{MobileNetV2 with vs.\ without fuzzy adaptive resampling.}
\label{tab:main}
\centering
\begin{tabular}{lccc}
\toprule
Metric & Without Fuzzy & With Fuzzy & $\Delta$ \\
\midrule
Overall Accuracy        & 63\% & 86\% & +23 \\
Macro-Avg F1            & 0.49 & 0.85 & +0.36 \\
Moderate Dem.\ Recall   & 0\%  & 90\% & +90 \\
Very Mild Dem.\ F1      & 0.17 & 0.88 & +0.71 \\
Mild Dem.\ Precision    & 0.06 & 0.89 & +0.83 \\
Non-Demented Specificity& 67\% & 96\% & +29 \\
\bottomrule
\end{tabular}
\end{table}

\subsection{Comparison with Standard Imbalance Methods}
\textbf{TODO (required for a journal):} report the fuzzy controller against
no-balancing, class weights, random oversampling, SMOTE, and focal loss on
identical folds. Without this table a reviewer cannot judge whether the fuzzy
step beats simpler alternatives.
\begin{table}[t]
\caption{Fuzzy resampling vs.\ standard imbalance baselines (TODO: fill).}
\label{tab:baselines}
\centering
\begin{tabular}{lcc}
\toprule
Method & Accuracy & Macro-F1 \\
\midrule
No balancing            & TODO & TODO \\
Class weights           & TODO & TODO \\
Random oversampling     & TODO & TODO \\
SMOTE                   & TODO & TODO \\
Focal loss              & TODO & TODO \\
\textbf{Fuzzy (ours)}   & TODO & TODO \\
\bottomrule
\end{tabular}
\end{table}

\section{Discussion and Limitations}
The adaptive controller substantially improves recall on rare stages and yields
clinically reasonable errors (confusions concentrate on adjacent stages).
However, several limitations temper the claims. First, the fuzzy policy is, in
the limit, close to proportional oversampling; the baseline comparison
(Table~\ref{tab:baselines}) is needed to establish that the fuzzy formulation
adds value. Second, results on the Kaggle dataset may be optimistic due to
potential subject-level leakage; external validation on an independent cohort
(e.g., ADNI/OASIS with patient-level splits) would strengthen the work. Third,
the membership functions are hand-set; sensitivity to their parameters should be
reported. Finally, this is a research/educational system and not a validated
diagnostic device.

\section{Conclusion}
We presented a lightweight, interpretable fuzzy controller for adaptive
class-rebalancing in Alzheimer's MRI staging, integrated with MobileNetV2. The
method recovers recall on rare dementia stages that a naive CNN misses. Pending
the baseline comparison and cross-validated, leakage-controlled evaluation, the
approach offers a simple and reproducible route to robust staging under class
imbalance.

\section*{Reproducibility}
Code and a live demo are available at
\url{https://github.com/shvm-k/AlzDetect}.

\begin{thebibliography}{10}
\small
% TODO: replace each entry with full author list, venue, year, pages, DOI.
\bibitem{ref1} ``Classification of Alzheimer's Disease Using Deep Learning based
Edge Detection and Fuzzy Neural Network,'' TODO.
\bibitem{ref2} ``FP-CNN: Fuzzy pooling-based convolutional neural network for
lung ultrasound image classification with explainable AI,'' TODO.
\bibitem{ref3} ``Fuzzy Deep Learning for the Diagnosis of Alzheimer's Disease:
Approaches and Challenges,'' TODO.
\bibitem{ref4} ``Integration of fuzzy logic and a convolutional neural network in
three-way decision-making,'' TODO.
\bibitem{ref5} ``Machine Learning Technique for Alzheimer's Disease (AD) --- A
Comprehensive Review,'' TODO.
\bibitem{ref6} ``Fuzzy Logic Module of Convolutional Neural Network for
Handwritten Digits Recognition,'' TODO.
\bibitem{ref7} ``A Convolutional Fuzzy Neural Network for Image Classification,''
TODO.
\bibitem{ref8} ``Review: Deep Learning and Fuzzy Logic Applications,'' TODO.
\bibitem{ref9} ``Transfer Learning based Image Visualization using CNN,'' TODO.
\bibitem{ref10} ``Convolutional Neural Network and Fuzzy Logic-based Hybrid
Melanoma Diagnosis System,'' TODO.
\end{thebibliography}

\end{document}
